\uuid{wkRy}
\titre{Ajuster une droite en pénalisant sa pente}
\niveau{L3}
\module{OPTRN}
\chapitre{Optimisation}
\sousChapitre{Régression et régularisation}
\theme{Moindres carrés, pénalisation}
\auteur{Maxime Nguyen}
\datecreate{2026-09-26}
\organisation{AMSCC}
\difficulte{2}

\contenu{
\texte{On veut ajuster le modèle $\widehat y=ax+b$ aux deux observations $(0,0)$ et $(1,1)$. Pour $\lambda\in\{0,1\}$, on minimise
\[
J_\lambda(a,b)=b^2+(a+b-1)^2+\lambda a^2.
\]
Le dernier terme pénalise les pentes de grande valeur.}

\begin{enumerate}
\item \question{Identifier dans $J_\lambda$ les erreurs commises sur les deux observations et le terme de pénalisation.}
\indication{Calculer la prédiction $ax+b$ en $x=0$ puis en $x=1$.}
\reponse{En $x=0$, la prédiction vaut $b$ et la cible vaut $0$ : le carré du résidu est $b^2$. En $x=1$, la prédiction vaut $a+b$ et la cible vaut $1$ : le carré du résidu est $(a+b-1)^2$. Le terme $\lambda a^2$ pénalise la pente.}

\item \question{Déterminer le minimum de $J_0$ et le modèle obtenu.}
\indication{La somme de deux carrés est positive ou nulle. Peut-on rendre les deux carrés nuls ?}
\reponse{Le minimum de $J_0=b^2+(a+b-1)^2$ vaut $0$. Il faut $b=0$ et $a+b=1$, d'où l'unique modèle $\widehat y=x$, soit $(a,b)=(1,0)$.}

\item \question{Calculer le gradient de $J_1$ et déterminer son unique minimiseur.}
\indication{Résoudre le système $\nabla J_1=0$, puis justifier l'unicité en examinant la hessienne.}
\reponse{On a
\[
\nabla J_1(a,b)=\begin{pmatrix}4a+2b-2\\2a+4b-2\end{pmatrix},
\qquad
\nabla^2J_1=\begin{pmatrix}4&2\\2&4\end{pmatrix}.
\]
Les valeurs propres de la hessienne sont $2$ et $6$ : elle est définie positive. Le système $2a+b=1$, $a+2b=1$ donne l'unique minimiseur $(a,b)=(1/3,1/3)$.}

\item \question{Comparer les pentes des deux modèles et leurs prédictions en $x=2$. Que change la pénalisation ?}
\indication{Utiliser les paramètres obtenus ; ne pas conclure à la qualité d'une prédiction sans connaître la vraie valeur en $x=2$.}
\reponse{Sans pénalisation, la pente vaut $1$ et la prédiction en $x=2$ vaut $2$. Avec $\lambda=1$, la pente vaut $1/3$ et la prédiction vaut $1$. La pénalisation réduit ici la pente, au prix d'erreurs non nulles sur les deux observations. On ne peut pas dire quel modèle prédit le mieux en $x=2$ sans connaître la valeur observée à cet endroit.}
\end{enumerate}
}
